\chapter{評価実験}
\label{sec:experiments}

\section{実験条件の例}
表\ref{tab:settings}は，比較条件の記載方法を示すために作成したものである．
実際にこの条件で実験を実施したわけではない．論文では，使用した設定と，その設定を選択した理由を示す．
結果を見た後で条件を変更した場合には，最初の計画との差異を説明し，追加の分析と区別する必要がある．

\begin{table}[htbp]
  \centering
  \caption{書式確認用の仮の比較条件}
  \label{tab:settings}
  \begin{tabular}{ll}\hline
    項目 & 設定例\\\hline
    候補文書数 & 各検索要求につき100件\\
    評価対象 & 上位10件の検索結果\\
    比較する構成 & 基準システムと並べ替えを行う構成\\
    集計単位 & 検索要求ごとの得点\\
    集計方法 & 各検索要求を等しい重みで平均\\\hline
  \end{tabular}
\end{table}

\section{評価値の集計}
検索要求$i$に対する評価値を$s_i$，比較対象の要求数を$n$とすると，平均値は次式で表される．
\begin{equation}
  \label{eq:mean}
  \bar{s}=\frac{1}{n}\sum_{i=1}^{n}s_i.
\end{equation}
式\eqref{eq:mean}は，各要求に同じ重みを与える集計である．
この式だけでは，適合性や順位をどのように評価値へ変換するかは定まらない．
実際には，$s_i$を求める指標の定義や，未判定文書の扱いも明示する必要がある．

実験の記録には，要求識別子と両システムの得点を残し，設定ファイルや実行ログと対応付ける．
この記録方法の改善案を次章で説明する．

% Demonstration-only break; flush the tables before showing a continuation page.
\clearpage
\section{結果と図の配置例}
表\ref{tab:results}の値は，桁の配置と参照の動作を確認するための仮の値である．
二つの群の要求数が等しいという設定の下で，全体の平均を示している．
この表から，実在のシステムの有効性や統計的な確かさを主張することはできない．

\begin{table}[htbp]
  \centering
  \caption{実験結果ではない仮の評価値}
  \label{tab:results}
  \begin{tabular}{lrrr}\hline
    検索要求の群 & 基準 & 比較対象 & 差\\\hline
    短い要求 & 0.420 & 0.460 & +0.040\\
    詳細な要求 & 0.380 & 0.390 & +0.010\\
    全体 & 0.400 & 0.425 & +0.025\\\hline
  \end{tabular}
\end{table}

図\ref{fig:logo}は，図の中央配置，図題，本文からの参照を確認するための例である．
元のテンプレートに含まれていた校徽を使用しており，研究結果を示すものではない．
実際の論文では，図から読み取るべき内容と，その内容が研究課題にどう関係するかを本文で説明する．

\insertfigure{waseda_logo.pdf}{書式確認用に掲載した早稲田大学の校徽}{fig:logo}{width=3cm}

\section{結果を説明する際の注意}
平均値の上昇は，すべての要求で改善したことを意味しない．
個別の得点を調べ，改善と悪化がどのように分布しているかを確認する必要がある．
大きく変化した要求を取り上げる場合には，選択した基準を記述し，印象的な例だけで全体の傾向を
説明しようとしないことが重要である．具体例は挙動を理解する助けになるが，発生頻度の推定とは異なる．

考察では，観察した事実と，その原因についての仮説を分ける．
例えば，ある要求で無関係な文書が上位に現れた場合，語の曖昧さが関係している可能性はある．
しかし，一つの結果だけで一般的な原因を確定することはできない．
追加の事例や条件を制御した実験によって，その解釈を支える証拠を集める必要がある．

